\section*{Background and Summary}
%%% https://www.nature.com/sdata/submission-guidelines
% - provide an overview of dataset, including the motivation for creating it, as well as outlining the potential reuse value where applicable. 
% - Any previous publications that used these data, in whole or in part, should be cited and briefly summarised. 
% - there is no formal mandate to cite a given amount of prior art for comparison, however we recommend that at least some other datasets or outputs relevant to the field are cited for readers' general interest.
% - do not include subjective claims on novelty, impact, or utility.


Systematic reviews synthesize evidence through protocolled searching, selection, and documentation of relevant literature, typically using explicit eligibility criteria and structured synthesis~\cite{gough:2012,liberati:2009,mulrow:1994}. Developed and institutionalized most prominently in biomedicine~\cite{chalmers:1993,guyatt:1992}, systematic review methodology has since become established across a range of scientific disciplines, alongside the broader expansion of evidence-based policy and practice agendas~\cite{simons:2021}. Methodological guides in biomedicine~\cite{chandler_cochrane_2019}, nursing~\cite{bettany_saltikov_how_2012}, education~\cite{newman_systematic_2020}, social science~\cite{petticrew_systematic_2008}, design~\cite{lame:2019}, and computer science~\cite{kitchenham:2004} reflect both shared review principles and field-specific adaptations in information sources, search conventions, eligibility criteria, and synthesis approaches. Across domains, reviews commonly proceed from defining objectives and eligibility criteria to retrieving and screening studies, extracting and appraising study data, and synthesizing findings. Because retrieval and screening are among the most time-consuming stages, substantial effort has focused on their automation. Evaluation resources for these tasks remain concentrated in biomedicine (see Table~\ref{table-sr-datasets}), leaving query formulations, search strategies, and relevance criteria from other fields sparsely represented.



% Systematic reviews have become a dominant method and, in many areas, the standard for synthesizing evidence to answer focused research questions, especially in medicine, education, and the social sciences. Systematic reviews follow a protocolled and transparent approach to searching for, selecting, and documenting relevant literature, with explicit inclusion and exclusion criteria. While they share this systematic approach to evidence identification and screening, they can differ substantially with respect to synthesis methods, ranging from qualitative evidence syntheses and mixed-methods reviews to quantitative meta-analyses that aggregate effect sizes~\cite{gough:2012}. Systematic reviews are valued for their completeness in capturing the available evidence, transparency in documenting procedures and data sources, reproducibility through well-defined methodologies, rigorous appraisal of study quality (including risk-of-bias assessment and related threats to validity), and their ability to support more robust and informative conclusions via structured synthesis across studies~\cite{liberati:2009, mulrow:1994}. Their results guide subsequent research, policy, and professional practice.

% Historically, systematic review methodology was developed and institutionalized most prominently in biomedicine to extract and communicate ``the evidence'' from an expanding primary literature for decision-making~\cite{chalmers:1993, guyatt:1992}. Related evidence-based policy and practice agendas then helped diffuse evidence-synthesis toolkits beyond health, supported by new institutions and infrastructures for producing and disseminating reviews~\cite{simons:2021}. Over time, these principles spread into many other disciplines, often with field-specific adaptations in information sources, search conventions, eligibility criteria, and synthesis approaches. Regardless of the domain and application, a systematic review is typically conducted in a multi-step workflow, from scoping the objective and inclusion criteria, through searching for and selecting studies (retrieval and screening), collecting and appraising study data, and finally synthesizing, presenting, and interpreting findings. Methodological guides in biomedicine~\cite{chandler_cochrane_2019}, nursing~\cite{bettany_saltikov_how_2012}, education~\cite{newman_systematic_2020}, social science~\cite{petticrew_systematic_2008}, design~\cite{lame:2019}, or computer science~\cite{kitchenham:2004} cover these core steps, while differing in granularity and how they group them into phases.

% Conducting a systematic review is time-consuming and expensive, and substantial effort has been devoted to automating one of its most expensive parts, namely retrieval and screening. Evaluation resources for retrieval and screening have, however, been largely concentrated in the biomedical domain (see Table~\ref{table-sr-datasets}), reflecting the empirical prominence of systematic reviews in biomedicine. As a result, domain-specific query formulations, search strategies, and relevance criteria outside biomedicine are sparsely represented in existing collections. This motivated the construction of a corpus that covers systematic reviews across scientific fields rather than a single domain.

Existing information retrieval test collections for evaluating retrieval and screening in systematic reviewing have predominantly focused on the biomedical domain, differing in scale, annotation depth, and supported tasks, but sharing the goal of providing reusable test collections for studying search, screening, and study prioritization. The SIGIR~2017 SysRev Query Collection~\cite{scells:2017} comprises Cochrane reviews published between~2014 and~2016, with the original Boolean search strategies translated into executable Elasticsearch queries and included and excluded studies mapped to PubMed identifiers. The Seed Studies Collection~\cite{wang:2022} targets seed-driven search strategies, providing original PubMed Boolean queries, the seed studies used to formulate them, lists of included studies, and references obtained through automated citation chasing tools such as CitationChaser and SpiderCite. The CLEF TAR collections~\cite{kanoulas:2017,kanoulas:2018,kanoulas:2019}, released as part of the CLEF lab on technology-assisted reviews in empirical medicine and built on Cochrane reviews, contain expert-formulated queries, retrieved records, and relevance assessments spanning intervention, prognosis, qualitative, and diagnostic test accuracy review types, primarily supporting querying and screening prioritization. Beyond these, \citet{alharbi:2019} provide a dataset for retrieval and screening on systematic review updates, where two versions of a review exist and the task is to identify relevant studies for the updated version, while the dataset of \citet{hannousse:2021}, used exclusively for screening, is one of the few non-biomedical systematic review collections. The CSMeD dataset~\cite{kusa:2023} consolidates nine prior biomedical and computer science screening datasets~\cite{cohen:2006,wallace:2010,howard:2016,scells:2017,kanoulas:2017,kanoulas:2018,kanoulas:2019,alharbi:2019,hannousse:2021} into a unified meta-collection of 325~systematic reviews with standardized access via the BigBio framework, and enriches a subset of Cochrane reviews with eligibility criteria and search strategies extracted from review protocols, along with CSMeD-FT, a full-text screening collection with time-based splits. The dataset of \citet{wang:2025b} was used primarily as training data for Boolean query formulation through reinforcement learning; it is the largest of these collections but provides fewer review-process artifacts. Across these resources, scope typically remains limited to a single scientific domain and a relatively small number of curated review topics.

\webisSRCorpus{} is a cross-disciplinary corpus of 334,893~systematic reviews spanning 27~scientific domains, derived from reviews indexed in OpenAlex~\cite{Priem2022OpenAlexAF}. A subset of reviews is linked to its resolved set of cited studies, providing stable reference-list information. In addition, a range of methodological information is extracted from the available review full texts where explicitly reported. For systematic reviews that report a search strategy, whether as explicit Boolean queries or keyword lists, a canonical Boolean representation is constructed as a minimal approximation of the reported strategy for execution in OpenAlex. 
The corpus is organized into layers comprising reference links, extracted method fields, and normalized queries, which can support a range of reuse scenarios across scientific domains, including retrieval and screening research, training and evaluation of extraction methods for systematic review artifacts, and comparative analyses of review practices across fields and over time.